\chapter{The Atom as a Resonant System and the Geometric Origin of Chemistry}

\author{James Tyndall}

%----------------------------------------------------------------------
\section{Introduction: From Particles to Patterns}
%----------------------------------------------------------------------

The most fundamental pattern in the universe is the atom.  The Standard Model explains atomic architecture through abstract, probabilistic quantum rules.  SDT provides a deterministic, mechanical alternative: an atom is a \textbf{stable, nested, three-dimensional resonant system}---displacement vortices held together and organised by pressure gradients and geometric constraints.


%----------------------------------------------------------------------
\section{The Atomic Pressure Field}
%----------------------------------------------------------------------

\subsection{The Nuclear Field}

The nucleus creates a powerful, propagating pressure field:
\begin{equation}\label{eq:P-nuc}
  P_{\text{prop}}(r) = P_{\text{surface}} \left(\frac{R_{\text{nucleus}}}{r}\right)^2
\end{equation}

\subsection{Resonant Cavities (Shells)}

The spinning, dynamic nuclear vortex forms nested, stable \textbf{standing-wave cavities} in the surrounding medium.  An electron vortex can only maintain stable resonance at these specific, quantised distances.  These cavities \emph{are} the electron shells.

\subsection{Shell Quantisation}

\begin{equation}\label{eq:shell-radii}
  r_n \approx a_0 \cdot \frac{n^2}{Z_{\text{eff}}}
\end{equation}

where $a_0$ is the Bohr radius, $n$ is the principal shell number, and $Z_{\text{eff}}$ is the effective nuclear charge accounting for inner-shell occlusion (``shielding'').  This is the SDT origin of the \textbf{periods} in the periodic table.


%----------------------------------------------------------------------
\section{Electron Vortex Geometry (Subshells)}
%----------------------------------------------------------------------

The ``blocks'' of the periodic table reflect the stable geometric shapes an electron vortex can form within a given shell.

\begin{center}
\begin{tabular}{clll}
\hline
$\ell$ & \textbf{Type} & \textbf{Geometry} & \textbf{First Appears} \\
\hline
0 & s-vortex & Spherically symmetric oscillation & $n = 1$ \\
1 & p-vortex & Axially-oriented, dumbbell-shaped & $n = 2$ \\
2 & d-vortex & Multi-lobed, five orientations & $n = 3$ \\
3 & f-vortex & Complex, seven orientations & $n = 4$ \\
\hline
\end{tabular}
\end{center}

\textbf{The Rule $\ell < n$} is a direct \textbf{geometric packing constraint}.  A complex, high-order vortex cannot form and maintain stable resonance within the tight, high-pressure confines of a low-$n$ shell.


%----------------------------------------------------------------------
\section{The Mechanical Origin of the Pauli Exclusion Principle}
%----------------------------------------------------------------------

In SDT, ``spin'' is the intrinsic \textbf{chirality} (handedness) of a vortex.

\textbf{The Geometric Rule:} A single resonant cavity can hold a maximum of two electron vortices, and \emph{only if they have opposite chirality}.

\textbf{The Mechanism:}
\begin{itemize}
  \item Two vortices of the \textbf{same chirality}: helical wakes interfere destructively $\to$ unstable.
  \item Two vortices of \textbf{opposite chirality}: can \textbf{geometrically mesh} like interlocking gears $\to$ stable spin-paired system with minimised external field.
\end{itemize}

This is the mechanical reality behind the Pauli Exclusion Principle.


%----------------------------------------------------------------------
\section{Chemical Properties: The Atomic k-Factor}
%----------------------------------------------------------------------

Chemical behaviour is encapsulated by the \textbf{atomic k-factor}:
\begin{align}
  v_{\text{valence}} &= \sqrt{2 E_{I1} / m_e} \label{eq:v-valence} \\
  k_{\text{atom}} &= c / v_{\text{valence}} \label{eq:k-atom}
\end{align}

\begin{center}
\begin{tabular}{crrrrl}
\hline
$Z$ & \textbf{Element} & $E_{I1}$ (eV) & $k_{\text{atom}}$ & $\chi_{\text{Pauling}}$ & \textbf{Character} \\
\hline
 3 & Li & 5.39 & 215.7 & 0.98 & Reactive alkali \\
 4 & Be & 9.32 & 172.9 & 1.57 & Alkaline earth \\
 5 & B  & 8.30 & 184.2 & 2.04 & Metalloid \\
 6 & C  & 11.26 & 156.4 & 2.55 & Nonmetal \\
 7 & N  & 14.53 & 138.8 & 3.04 & Nonmetal \\
 8 & O  & 13.62 & 142.1 & 3.44 & Nonmetal \\
 9 & F  & 17.42 & 128.0 & 3.98 & Halogen \\
10 & Ne & 21.56 & 113.6 & --- & Inert \\
\hline
\end{tabular}
\end{center}

\textbf{Trends:} $k_{\text{atom}}$ is highest for alkali metals (loosely bound, ``fluffy'' displacement field) and decreases across each period as increasing nuclear charge tightens the resonance.  High $k_{\text{atom}}$ means easy electron release---metallic character.


%----------------------------------------------------------------------
\section{Conclusion}
%----------------------------------------------------------------------

The periodic table is a \textbf{map of stable geometric solutions for packing resonant displacement vortices}:
\begin{itemize}
  \item \textbf{Shells ($n$)}: stable resonant cavities.
  \item \textbf{Subshells ($\ell$)}: allowed geometric shapes.
  \item \textbf{Pauli Principle}: geometric compatibility for chiral pairing.
  \item \textbf{Chemical properties}: consequences of the derived atomic k-factor.
\end{itemize}

The laws of chemistry are the emergent laws of the geometry of stable, resonant patterns in a pressurised medium.
